%%%%%%%%%%%%%%%%%%%%%%%%%%%%%%%%%%%%%%%%%%%%%%%%%%%%%%%%%%%%%%%%%%%%%%%%%%
%   FEHIME CEREN AY - CV
%   Tilpasset: Dataanalytiker, DNB Bank ASA (FINN-kode 468593513)
%%%%%%%%%%%%%%%%%%%%%%%%%%%%%%%%%%%%%%%%%%%%%%%%%%%%%%%%%%%%%%%%%%%%%%%%%%

\documentclass[10pt,a4paper,roman]{moderncv}

\moderncvstyle{casual}
\moderncvcolor{blue}

\usepackage[utf8]{inputenc}
\usepackage[norsk]{babel}
\usepackage[scale=0.78,top=1.6cm,bottom=1.4cm]{geometry}
\setlength{\hintscolumnwidth}{3.4cm}
\usepackage{enumitem}
\setlist[itemize]{leftmargin=*,topsep=2pt,itemsep=1pt,parsep=0pt}

% Kompakt punktliste til bruk i cvitem
\newcommand{\cvbullets}[1]{\begin{itemize}#1\end{itemize}}

%----------------------------------------------------------------------------------------
%   NAVN OG KONTAKT
%----------------------------------------------------------------------------------------

\firstname{Fehime Ceren}
\familyname{Ay}

\title{Dataanalytiker \textbar{} PhD i økonomi}
\address{Bestumveien 37C, 0281}{Oslo}
\phone{+47 92 25 66 45}
\email{fcerenay@gmail.com}
\extrainfo{\href{https://cerenay.github.io}{\color{blue}cerenay.github.io}}

\begin{document}
\makecvtitle

%----------------------------------------------------------------------------------------
%   PROFIL
%----------------------------------------------------------------------------------------

\section{Profil}
\cvitem{}{
Dataanalytiker med PhD i økonomi og fem års erfaring fra analysearbeid i to store,
datadrevne organisasjoner: Skatteetaten og Telenor. Jeg jobber i hele analyseløpet:
fra å oversette uklare forretningsbehov til konkrete analysespørsmål, via uttrekk og
kvalitetssikring av data, til analyse og formidling av funn til både tekniske og
ikke-tekniske beslutningstakere. Til daglig jobber jeg i SQL, Python og R på
Databricks. Jeg er vant til å jobbe
selvstendig i problemstillinger som ikke er ferdig definerte, og til å samarbeide tett
med data engineers, IT- og AI-miljøer om felles datagrunnlag, begrepsdefinisjoner og
gjenbrukbare dataprodukter. Motivasjonen min er å sette kunnskap i system slik at
beslutninger tas på et bedre grunnlag, ikke på magefølelse.
}

%----------------------------------------------------------------------------------------
%   NØKKELKOMPETANSE
%----------------------------------------------------------------------------------------

\section{Nøkkelkompetanse}

\cvitem{Programmering}{SQL (uttrekk, modellering, optimalisering), Python, R}
\cvitem{Dashbord \& formidling}{Power BI, R Shiny og Quarto til utvikling av
selvbetjente analyseverktøy og rapporteringsløsninger}
\cvitem{Plattform \& data}{Databricks, skybaserte og lokale datamiljøer,
Git/versjonskontroll, arbeid med store register- og loggdatasett}
\cvitem{Metode}{Økonometri og kausal inferens, kvasi-eksperimentelle metoder,
eksperimentdesign og A/B-testing, prediktiv modellering, spørreundersøkelser
og effektmåling}
\cvitem{Datakvalitet}{Datavalidering og avviksanalyse, dokumentasjon av
definisjoner og måltall, etablering av felles datagrunnlag på tvers av kilder}
\cvitem{Språk}{Norsk (flytende), Engelsk (flytende), Tyrkisk (morsmål)}
\cvitem{Øvrig}{Norsk og tyrkisk statsborgerskap. Referanser oppgis på forespørsel.}

%----------------------------------------------------------------------------------------
%   ARBEIDSERFARING
%----------------------------------------------------------------------------------------

\section{Arbeidserfaring}

\cvitem{09/2023 -- nå}{\textbf{Seniorrådgiver / Dataanalytiker (atferd og effekt)} \newline
\textit{Skatteetaten, Skattedirektoratet -- Data, Analyse og Kunnskap}}
\cvitem{}{\vspace{-4pt}
\cvbullets{
\item Gjennomfører analyser og datauttrekk som gir beslutningsstøtte til
      fagavdelinger og ledelse, basert på store registerdatasett, loggdata og
      surveydata i Databricks.
\item Oversetter uklare forretningsbehov til konkrete analyseproblemstillinger,
      definerer databehov og leverer anbefalinger som faktisk tas i bruk.
\item Designer og analyserer eksperimenter, A/B-tester og effektmålinger som
      dokumenterer hvilke tiltak som virker, og hvilke som ikke gjør det.
\item Bygger analyseverktøy og rapporter (R Shiny, Power BI, Quarto) slik at
      fagmiljøene selv kan følge opp måltall uten å være avhengige av én person.
\item Arbeider med datakvalitet og begrepsforvaltning: avdekker avvik mellom
      kilder og bidrar til omforente definisjoner av sentrale måltall.
\item Samarbeider tett med IT, analyse- og AI-miljøer om datadrevne løsninger;
      deltar i internasjonale fagmiljøer (OECD-fokusgrupper, IBPPA).
}}

\cvitem{01/2021 -- 09/2023}{\textbf{Research Scientist / Dataanalytiker} \newline
\textit{Telenor ASA -- Research \& Innovation, Advanced Analytics \& AI}}
\cvitem{}{\vspace{-4pt}
\cvbullets{
\item Analyserte store interne datakilder (bl.a. kundeservice- og chatbot-logger,
      bruks- og aktivitetsdata) og koblet dem mot eksterne datakilder for å svare
      på konkrete forretningsspørsmål.
\item Utviklet prediktive modeller og NLP-baserte analyser av samtalelogger for å
      måle og forbedre kundetilfredshet. Arbeidet er publisert i
      \textit{Quality and User Experience}.
\item Strukturerte og kvalitetssikret rådata til analyseklare datasett, og
      dokumenterte definisjoner slik at analysene kunne gjenbrukes av andre team.
\item Leverte innsikt og anbefalinger på tvers av forretningsenheter i Norden og
      Asia, til både produktledelse, teknologimiljøer og konsernledelse.
\item Tilknyttet forsker ved NHH, FAIR Research Centre; ledet eksterne samarbeid
      med akademiske og profesjonelle miljøer.
}}

\cvitem{08/2016 -- 12/2020}{\textbf{PhD-stipendiat} \newline
\textit{Norges Handelshøyskole (NHH) -- Institutt for økonomi, FAIR / The Choice Lab}}
\cvitem{}{\vspace{-4pt}
Kvantitativ forskning på økonomisk beslutningstaking: eksperimentdesign,
datainnsamling og økonometrisk analyse av store forsøksdatasett, i samarbeid med
internasjonale forskergrupper. Underviste og veiledet i statistikk og
forskningsmetode.
}

\cvitem{10/2014 -- 08/2016}{\textbf{Forskningsassistent} \newline
\textit{Istanbul Universitet -- Fakultet for økonomi}}
\cvitem{}{\vspace{-4pt}
Datainnsamling og kvantitativ analyse i forskningsprosjekter innen offentlig
økonomi og skatteatferd.
}

%----------------------------------------------------------------------------------------
%   UTDANNING
%----------------------------------------------------------------------------------------

\section{Utdanning}

\cvitem{2016 -- 2021}{\textbf{PhD i økonomi} \newline
\textit{Norges Handelshøyskole (NHH) -- FAIR, The Choice Lab} \newline
Avhandling: \textit{Essays on Information Preferences and Morality} \newline
Veiledere: Erik Ø. Sørensen (NHH) og Pedro Dal Bó (Brown University) \newline
Forskningsopphold ved Brown University (2019)}

\cvitem{2014 -- 2016}{\textbf{Mastergrad i økonomi} \newline
\textit{Galatasaray Universitet, Tyrkia} \newline
Masteroppgave: \textit{A Voting Experiment on Promises and Fairness Perceptions}}

\cvitem{2009 -- 2014}{\textbf{Bachelor i økonomi og offentlig finans} \newline
\textit{Hacettepe Universitet, Tyrkia} \newline
Bacheloroppgave: \textit{Game Theoretical Approach to Tax Evasion and the Vickrey--Clarke--Groves Mechanism}}

%----------------------------------------------------------------------------------------
%   PUBLIKASJONER
%----------------------------------------------------------------------------------------

\section{Publikasjoner og analyser (utvalg)}

\cvitem{2026}{Hetland, A. G., Ay, F. C., Serdarevic, N., \& Lokken, N. C.
\textit{Noen Airbnb-utleiere rapporterer ikke leieinntektene sine -- hvorfor?}
Skatteetatens analysenytt.
\href{https://www.skatteetaten.no/om-skatteetaten/analyse-og-rapporter/analysenytt/airbnb-utleiere-rapporter-ikke-leieinntektene-sine/}{\color{blue}\mbox{Skatteetatens Analysenytt}}}

\cvitem{2025}{Ay, F. C., Freddi, E., Følstad, A., Hodnebrog, S., Kvale, K., Sell, O. A., \& Ulsaker, S.
\textit{Conversation logs as a source of insight: Predicting user satisfaction for customer service chatbots.}
Quality and User Experience, 10(1).}

\cvitem{2024}{Ay, F. C., Berge, J. W., \& Nødtvedt, K. B.
\textit{Strategic curiosity: An experimental study of curiosity and dishonesty.}
Journal of Economic Behavior \& Organization, 217, 287--297.}

\cvitem{2023}{Huber, C., Dreber, A., Huber, J., \textit{et al.} (inkl. Ay, F. C.)
\textit{Competition and moral behavior: A meta-analysis of forty-five crowd-sourced experimental designs.}
PNAS, 120(23).}

\cvitem{2023}{Azevedo, F., Pavlovic, T., Rigo, G. G., Ay, F. C., \textit{et al.}
\textit{Social and moral psychology of COVID-19 across 69 countries.}
Nature Scientific Data, 10(1), 272.}

\cvitem{2022}{Van Bavel, J. J., Cichocka, A., Capraro, V., \textit{et al.} (inkl. Ay, F. C.)
\textit{National identity predicts public health support during a global pandemic.}
Nature Communications, 13(1), 517.}

\cvitem{}{Fullstendig publikasjonsliste: \href{https://cerenay.github.io}{cerenay.github.io}}

\end{document}
